\documentclass[11pt]{article}
\usepackage[utf8]{inputenc}
\usepackage[T1]{fontenc}
\usepackage{lmodern}
\usepackage[margin=1in]{geometry}
\usepackage{amsmath,amssymb,mathtools}
\usepackage{graphicx}
\usepackage{grffile}
\usepackage{float}
\usepackage{hyperref}
\usepackage{parskip}
\usepackage{enumitem}
\usepackage{xurl}
\setlength{\parskip}{0.65em}
\setlength{\parindent}{0pt}
\begin{document}
\title{Daily Research Digest --- 2026-04-15}
\author{}
\date{}
\maketitle
\textbf{Topics:} galaxy clusters, galaxies, gravitational lensing, dark matter.
\section*{Synthesis}
\_Digest generation failed: timed out\_
\newpage
\section*{Paper Summaries and Figures}
\subsection*{1. Revisiting the angular size-redshift cosmological test with milliarcsecond radio structures in active galactic nuclei}
\textbf{Focus:} dark matter\\
\textbf{Authors:} Mina Ghodsi Yengejeh, Tatiana A. Koryukova, Leonid I. Gurvits, Sándor Frey, Alexander B. Pushkarev, Alexander V. Plavin, Kenneth I. Kellermann, András Kovács\\
\textbf{Published:} 2026-04-14T16:26:57+00:00\\
\textbf{PDF:} \href{https://arxiv.org/pdf/2604.12936v1}{https://arxiv.org/pdf/2604.12936v1}\\
\textbf{Summary:}
VLBI measurements of the sizes of compact extragalactic radio sources, jetted active galactic nuclei, provide data for probing the angular size--redshift relation, offering a complementary test to other distance--redshift methods. We analyse a significantly expanded dataset to reassess the angular size--redshift relation and its potential for constraining cosmological model parameters, focusing on the matter density parameter $Ω_{\mathrm{m}}$ in a flat $Λ$ Cold Dark Matter Universe. This is the first major update of the compact-source angular size test in the past quarter of a century, using a dataset an order of magnitude larger than in previous studies. MCMC analysis on real data and on multiple mock catalogues to evaluate parameter constraints in the presence of observational scatter. In addition, we conducted a test with 100 randomized catalogues created by shuffling redshifts while preserving other observables. We also explored how astrophysical parameters depend on fixed cosmological models with different $Ω_{\mathrm{m}}$ values. The randomization test showed that the posterior distributions from randomized data do not overlap with those from real observations, confirming that the measured angular size--redshift relation is physically meaningful. The astrophysical model parameter that describes the redshift dependence of the source angular size exhibits degeneracy with $Ω_{\mathrm{m}}$. Simulated mock catalogues indicate that the method is able to constrain $Ω_{\mathrm{m}}$ if the data scatter is below $\sim20\%$. Scaling estimates suggest that high-quality data of samples of several thousands to $\sim 100\,000$ sources, a standardisation calibration approach, and/or refining sample selection criteria are needed to fully exploit the potential of the angular size--redshift test with this type of objects (abridged).
\newpage
\subsection*{2. Infrared Spectral Gap in a Gluonic Dark Sector as the Origin of the Galactic Acceleration Scale}
\textbf{Focus:} dark matter\\
\textbf{Authors:} Gilles Cohen-Tannoudji, Jean-Pierre Gazeau, Hamed Pejhan, Jean-Pierre Treuil\\
\textbf{Published:} 2026-04-14T15:58:18+00:00\\
\textbf{PDF:} \href{https://arxiv.org/pdf/2604.12910v1}{https://arxiv.org/pdf/2604.12910v1}\\
\textbf{Summary:}
The radial acceleration relation reveals a nearly universal acceleration scale of order $10^{-10}\,\mathrm{ms^{-2}}$ in galactic dynamics, whose origin remains unexplained within conventional cold dark matter scenarios. We propose that this scale arises from an intrinsic infrared spectral property of the dark sector. Specifically, we hypothesize that a long-lived, color-neutral gluonic vacuum component survives the post-inflationary expansion era and, at large distances, develops a spectrally rigid lowest-weight structure. The microscopic seed for this infrared organization is provided by the QCD trace anomaly, which breaks classical scale invariance and, through dimensional transmutation, generates an intrinsic infrared scale in the gluonic sector. Requiring Lorentz covariance together with a positive-energy lowest-weight unitary realization then selects the Anti de Sitter algebra $\mathfrak{so}(2,3)$ as the simplest symmetry admitting a discrete tower of states with a representation-theoretically protected gap. The associated gap introduces a finite correlation length $r_{\texttt{c}}$ that controls the large-scale coherence of the dark sector. A self-gravitating condensate dominated by the lowest-weight mode generates a characteristic acceleration $g^{}_\star = G M_h / r_{\texttt{c}}^2$, naturally of the same order as the observed galactic acceleration scale, within standard Newtonian gravity. In this framework, the galactic acceleration scale appears as the gravitational imprint of a trace-anomaly-seeded infrared spectral gap in a coherent gluonic dark sector, rather than as a consequence of modified gravity or of galaxy-by-galaxy formation histories.
\subsubsection*{Selected figures}
\begin{figure}[H]
\centering
\includegraphics[width=\linewidth,height=0.42\textheight,keepaspectratio]{/home/kf/research-assistant/data/cache/figures/http-arxiv.org-abs-2604.12910v1/figure\_1.png}
\end{figure}
\newpage
\subsection*{3. Hot blue progenitors of stellar-mass black holes}
\textbf{Focus:} galaxies\\
\textbf{Authors:} Avishai Gilkis, Eva Laplace, Maria Drout, Charles Kilpatrick, Anna O'Grady, Christopher Tout\\
\textbf{Published:} 2026-04-14T15:18:12+00:00\\
\textbf{PDF:} \href{https://arxiv.org/pdf/2604.12868v1}{https://arxiv.org/pdf/2604.12868v1}\\
\textbf{Summary:}
While the connection between massive stars and supernova explosions is well established observationally, the link between massive stars and black hole formation remains elusive. Some massive stars may collapse directly to black holes without a successful supernova, and may therefore appear as disappearing stars. We investigate the expected photometric properties of such black hole progenitors by combining detailed single and binary stellar evolution models with physically motivated prescriptions linking pre-collapse core structure to explosion or direct collapse outcome, together with stellar atmosphere calculations, producing synthetic photometry across standard ultraviolet to infrared filters. Weighting by an initial mass function and empirical binary distributions, we predict both the observable distribution of progenitor brightness and colour and the rate of direct-collapse events, which we estimate to be about 0.4 per century for a galaxy forming stars at 1 Msun/yr. We find that black hole progenitors are predominantly hot and blue at the pre-collapse stage, with many in Wolf-Rayet phases and luminous in the ultraviolet, while only a minority are red supergiants. Consequently, searches that focus primarily on red supergiants are likely to miss a substantial fraction of direct-collapse events. Monitoring campaigns that include ultraviolet-sensitive observations of nearby star-forming galaxies therefore provide a promising route to detecting disappearing massive stars, offering a direct observational probe of black hole formation. Our results provide predictions to interpret such surveys and constrain the channels that lead to black hole formation.
\subsubsection*{Selected figures}
\begin{figure}[H]
\centering
\includegraphics[width=\linewidth,height=0.42\textheight,keepaspectratio]{/home/kf/research-assistant/data/cache/figures/http-arxiv.org-abs-2604.12868v1/figure\_1.png}
\end{figure}
\newpage
\subsection*{4. Dark Matter's influence on Evolution of MBHB in Dwarf Galaxies: A Case Study of Leo I dSph}
\textbf{Focus:} dark matter\\
\textbf{Authors:} Muhammad Junaid, Fazeel Mahmood Khan\\
\textbf{Published:} 2026-04-14T12:25:59+00:00\\
\textbf{PDF:} \href{https://arxiv.org/pdf/2604.12658v1}{https://arxiv.org/pdf/2604.12658v1}\\
\textbf{Summary:}
In this study, we investigate the dynamical evolution of a massive binary black hole (MBHB) in the Leo I dwarf spheroidal galaxy model and examine how dark matter along with stellar matter's gravitational interactions influence its long-term behavior. Using high-resolution direct N-body simulations, we follow the orbital evolution of the binary within a realistic model of the Leo I stellar and dark matter distribution. We found that the binary separation decreases from an initial 300-parsec orbit to roughly 1 parsec over a period of about 2 Gyr, primarily driven by dynamical friction and stellar hardening. The orbital evolution then stalls at this scale, illustrating the well-known final parsec problem. During this phase, the binary also develops increasing orbital eccentricity and produces a modest redistribution of the inner mass profiles in some cases. We then further estimate the final stage of the system's evolution using gravitational-wave emission models and find that the binary is unlikely to merge within a Hubble time. The prolonged dynamical friction phase appears to be related to the low stellar and dark matter densities in Leo I. These results suggest that massive binary black holes in dwarf spheroidal galaxies such as Leo I will not contribute to the gravitational-waves detectable from LISA even if dark matter is considered.
\subsubsection*{Selected figures}
\begin{figure}[H]
\centering
\includegraphics[width=\linewidth,height=0.42\textheight,keepaspectratio]{/home/kf/research-assistant/data/cache/figures/http-arxiv.org-abs-2604.12658v1/figure\_1.png}
\end{figure}
\newpage
\subsection*{5. Primordial Black Holes Formation Beyond the Standard Cosmic QCD Transition}
\textbf{Focus:} dark matter\\
\textbf{Authors:} Maël Gonin, Oleksii Ivanytskyi, David Blaschke, Günther Hasinger\\
\textbf{Published:} 2026-04-14T11:04:50+00:00\\
\textbf{PDF:} \href{https://arxiv.org/pdf/2604.12581v1}{https://arxiv.org/pdf/2604.12581v1}\\
\textbf{Summary:}
We review the role of primordial black holes (PBHs) for illuminating the dark ages of the cosmological evolution and as dark matter (DM) candidates. We elucidate the role of phase transitions for primordial black hole formation in the early Universe and focus our attention on the cosmological QCD phase transition within a recent microscopical model. We explore the impact of physics beyond the Standard Model (SM) on the cosmic equation of state and the probability distribution for the formation of PBHs which serve as candidates for DM and contribute to present-day binary black-hole merger events.
\subsubsection*{Selected figures}
\begin{figure}[H]
\centering
\includegraphics[width=\linewidth,height=0.42\textheight,keepaspectratio]{/home/kf/research-assistant/data/cache/figures/http-arxiv.org-abs-2604.12581v1/figure\_1.png}
\end{figure}
\newpage
\subsection*{6. Two component pseudo-Nambu-Goldstone-boson dark matter}
\textbf{Focus:} dark matter\\
\textbf{Authors:} Riasat Sheikh, Takashi Toma, Koji Tsumura\\
\textbf{Published:} 2026-04-14T08:26:55+00:00\\
\textbf{PDF:} \href{https://arxiv.org/pdf/2604.12435v1}{https://arxiv.org/pdf/2604.12435v1}\\
\textbf{Summary:}
We study a two-component pseudo-Nambu-Goldstone-boson (pNGB) dark matter (DM) model motivated by boosted dark matter (BDM). The model is based on a complex scalar field charged under a dark $\text{U}(1)_V$ gauge symmetry, with a softly broken global $\text{SU}(3)_g$ symmetry that is spontaneously broken. The pNGB nature suppresses DM--Nucleon scattering, while the residual $\text{U}(1)_3 \times \text{U}(1)_{T_0}$ symmetry automatically stabilizes the two pNGB DM candidates and allows conversion of the heavier component into the lighter one. A central point is that the heavier or light component hierarchy is controlled by the two independent soft-breaking parameters that split the pNGB multiplet, so an abundant heavier component required for BDM can be obtained without introducing ad hoc hierarchies among independent portal coupling tuned to enable effective conversion. We analyze the relic abundance together with the constraints considered in this work, including Higgs invisible decays and perturbative unitarity, classify the coupled freeze-out dynamics, and assess the resulting BDM scattering cross section and flux.
\subsubsection*{Selected figures}
\begin{figure}[H]
\centering
\includegraphics[width=\linewidth,height=0.42\textheight,keepaspectratio]{/home/kf/research-assistant/data/cache/figures/http-arxiv.org-abs-2604.12435v1/figure\_1.png}
\end{figure}
\newpage
\subsection*{7. Probing Collapsed Dark Matter Halos with Fast Radio Bursts}
\textbf{Focus:} gravitational lensing, dark matter\\
\textbf{Authors:} Yuxuan He, Weiyang Wang, Chen Zhang, Yi-Ming Zhong\\
\textbf{Published:} 2026-04-14T01:41:46+00:00\\
\textbf{PDF:} \href{https://arxiv.org/pdf/2604.12189v1}{https://arxiv.org/pdf/2604.12189v1}\\
\textbf{Summary:}
Observations of ultra-dense substructures in strong lensing systems challenge the standard cosmological model at small scales. Self-interacting dark matter (SIDM), as an alternative to the cold and collisionless dark matter (CDM) of the standard cosmological model, provides a natural mechanism for forming such structures via gravothermal core collapse. We show that strong gravitational lensing of fast radio bursts (FRBs) provides an effective approach to detecting these substructures and probing dark matter self-interactions. Core-collapsed SIDM halos exhibit steeper central density profiles than CDM halos, enhancing the lensing cross section and producing longer time delays between FRB images. We compute lensing properties of core-collapsed subhalos and host halos, including maximal impact parameters and time-delay distributions. We demonstrate that future all-sky monitors, such as BURSTT, SKA2-Low, and SKA2-Mid, which are expected to detect $10^{5}$--$10^{7}$ FRBs over a decade, can measure time-delay distributions with high statistical significance. Modeling collapsed halos with a cored power-law density profile with inner slope $\gamma =3$ and assuming no excess beyond the singular isothermal sphere lens model, we show that our strategy can probe self-interaction cross section strengths of $\sigma _{\text{SI}}/m \gtrsim \min\{18,\, 40\lambda _{\text{sub}}\}\,\text{cm}^2/\text{g}$, where $\lambda _{\text{sub}}$ parameterizes the collapse time of a subhalo relative to that of the isolated case.
\newpage
\subsection*{8. The 3-3-1 Model: a natural framework for sub-MeV dark matter}
\textbf{Focus:} dark matter\\
\textbf{Authors:} Vinicius Oliveira, D. Cogollo, A. Doff, C. A. de S. Pires\\
\textbf{Published:} 2026-04-13T22:56:01+00:00\\
\textbf{PDF:} \href{https://arxiv.org/pdf/2604.12121v1}{https://arxiv.org/pdf/2604.12121v1}\\
\textbf{Summary:}
We show that the $\mathrm{SU}(3)_C \times \mathrm{SU}(3)_L \times \mathrm{U}(1)_N$ model with right-handed neutrinos naturally accommodates a viable sub-MeV dark matter (DM) candidate realized as pseudo-Goldstone boson that acquires tiny mass through gravitational effects. The observed relic abundance is obtained via freeze-in in a low-reheating temperature scenario, without requiring tiny couplings. The model operates at the TeV scale and remains testable at current and future collider experiments.
\newpage
\subsection*{9. Distinguish Bardeen-like black holes by Gravitational lensing}
\textbf{Focus:} gravitational lensing\\
\textbf{Authors:} Limei Yuan, Chen-Hung Hsiao, Yidun Wan\\
\textbf{Published:} 2026-04-13T18:44:16+00:00\\
\textbf{PDF:} \href{https://arxiv.org/pdf/2604.11951v1}{https://arxiv.org/pdf/2604.11951v1}\\
\textbf{Summary:}
We study Bardeen-like regular black holes without Cauchy horizons via gravitational lensing. In the weak field, the deflection angle receives a positive $\ell$-dependent correction, producing a slightly larger Einstein ring. For the galaxy ESO 325-G004, the predicted ring radius is consistent with current observations. In the strong field, for Sgr A* and M87*, the asymptotic position $θ_{\infty}$ remains identical to the Schwarzschild value; however, SDL coefficients are $\ell$-dependent, the angular separation s increases and the relative flux ratio $r_{\mathrm{mag}}$ decreases as $\ell$ increases. Time delays between relativistic images for Sgr A* and M87* also increase mildly with $\ell$. Our calculated values for these observables remain consistent with current observations. Future strong-field measurements of $\Delta T_{2,1}$, s, and $r_{\mathrm{mag}}$ may offer a viable test for regular black holes free of Cauchy horizons and may distinguish Bardeen-like from Schwarzschild black holes.
\subsubsection*{Selected figures}
\begin{figure}[H]
\centering
\includegraphics[width=\linewidth,height=0.42\textheight,keepaspectratio]{/home/kf/research-assistant/data/cache/figures/http-arxiv.org-abs-2604.11951v1/figure\_1.png}
\end{figure}
\newpage
\subsection*{10. Gravitational lensing around a Kerr-Sen black hole in plasma background}
\textbf{Focus:} gravitational lensing\\
\textbf{Authors:} Saswati Roy, Shubham Kala, Sayanika Modak, Hemwati Nandan, Amare Abebe\\
\textbf{Published:} 2026-04-13T12:30:21+00:00\\
\textbf{PDF:} \href{https://arxiv.org/pdf/2604.11860v1}{https://arxiv.org/pdf/2604.11860v1}\\
\textbf{Summary:}
We investigate the gravitational lensing of massless particles around a Kerr-Sen black hole immersed in a magnetized, cold, pressureless plasma medium. Both homogeneous and inhomogeneous plasma distributions are considered in this study to mimic realistic astrophysical environments. The light deflection angle is computed, and the effects of the black hole's rotation and charge on light bending are analyzed in detail. The conditions for the circular photon orbits are also examined in both plasma configurations. A comparison with the vacuum case (i.e. zero plasma frequency) highlights the role of plasma in modifying the light propagation, and the results obtained provide a deeper insight into plasma effects which improve our understanding of observational signatures of rotating charged black holes.
\subsubsection*{Selected figures}
\begin{figure}[H]
\centering
\includegraphics[width=\linewidth,height=0.42\textheight,keepaspectratio]{/home/kf/research-assistant/data/cache/figures/http-arxiv.org-abs-2604.11860v1/figure\_1.png}
\end{figure}
\end{document}
